\section{Representation by schemes}

\begin{lemma}[Hom-families are bounded]
\label{akcp-6-1-hom-family}
\dcref{6.1}
\source{MR0555258-compactifying-picard:page-0041}
If $X/S$ is strongly projective and two families of coherent fiber sheaves are
$b$-families (respectively $m$-regular families) with finitely many Hilbert
polynomials, then the family of their Hom sheaves is again a bounded,
$b$-family and $m$-regular family with finitely many polynomials.
\end{lemma}
\begin{proof}
Embed $X$ in $\PP(E)$.  Limitedness realizes both families as fibers of a
single finite-type parameter scheme after flattening.  A finite locally free
presentation of the first universal module expresses Hom as the kernel of a
map between coherent modules.  Flatten the cokernel so formation of this kernel
commutes with base change.  The bounded-family and regularity theorems for
kernels give the assertion.
\end{proof}

\begin{proposition}[Parameter scheme for the compactified Picard functor]
\label{akcp-6-2-parameter}
\dcref{6.2}
\source{MR0555258-compactifying-picard:page-0042}
\uses{akcp-1-3-local-free,akcp-2-6-quot-representability,akcp-2-9-effective-quotient,akcp-3-4-hom-rank-one,akcp-5-17-abel-fiber,akcp-6-1-hom-family}
Let $X/S$ be flat, finitely presented and projective with integral fibers of
dimension $r$, with a fixed very ample sheaf and one Hilbert polynomial.  Let
$\FF$ be a relatively rank-one torsion-free sheaf of one polynomial $\psi$.
For a polynomial $\epsilon$, let $P$ be the etale subsheaf of the Picard
functor consisting of classes $[\II]$ satisfying, on every fiber,
\[
\chi(\II(n))=\chi(\FF(n))-\epsilon(n),\qquad
\Ext^1(\II,\FF)=0.
\]
Then $P$ is represented by a strongly quasi-projective $S$-scheme.
\end{proposition}
\begin{proof}
Choose a uniform regularity bound for the family of possible $\II$, for $\FF$,
and for their Hom sheaves using \cref{akcp-3-4-hom-rank-one} and
\cref{akcp-6-1-hom-family}.
Replace $\FF$ by a twist so these families are regular from degree zero.  The
inverse image of $P$ under the Abel map is an open retrocompact subscheme
$Z$ of $Q=\Quot^\epsilon_{\FF/X/S}$: the pseudo-ideal of the universal quotient
is torsion-free rank one, and the Ext vanishing is open by
\cref{akcp-1-10-open-ext}.

For a point represented by $\II$, the fiber is
$\PP(H(\II,\FF))$ by \cref{akcp-5-17-abel-fiber}.  The local-to-global Ext
spectral sequence and regularity imply $H(\II,\FF)$ is locally free and
nonzero, so these projective fibers are nonempty.  Uniform regularity and
determinants of the exact sequence
$0\to\II(m)\to\FF(m)\to\GG(m)\to0$ show that their degrees are constant.
Thus the induced equivalence relation $\PP(H)\rightrightarrows Z$ is flat,
proper, finitely presented and has finitely many Hilbert polynomials.  The
quotient theorem \cref{akcp-2-9-effective-quotient} produces a strongly
quasi-projective scheme $P'$.  Its quotient map is an etale epimorphism, and
both $P'$ and the original etale sheaf are the quotient of the same relation;
therefore they represent the same functor.
\end{proof}

\begin{theorem}[Picard components for integral fibers]
\label{akcp-6-3-picard-integral}
\dcref{6.3}
\source{MR0555258-compactifying-picard:page-0045}
\uses{akcp-5-13-picard-open,akcp-6-2-parameter}
If $X/S$ is flat projective with geometrically integral fibers of fixed
dimension and fixed polarization polynomial, every polynomial component
$\Pic^{\varphi,\mathrm{tf}}_{X/S}$ is strongly quasi-projective.
\end{theorem}
\begin{proof}
Apply \cref{akcp-6-2-parameter} with $\FF=\OO_X$; the torsion-free Picard
component is an open retrocompact subsheaf by \cref{akcp-5-13-picard-open}.
\end{proof}

\begin{corollary}[Locally projective Picard functor]
\label{akcp-6-4-picard-locally-projective}
\dcref{6.4}
\source{MR0555258-compactifying-picard:page-0046}
\uses{akcp-6-3-picard-integral,akcp-5-8-spl-polynomial-open}
For flat finitely presented projective $X/S$ with geometrically integral fibers,
$\Pic^{\mathrm{tf},\acute{\mathrm et}}_{X/S}$ is represented by a disjoint union
of quasi-projective $S$-schemes.
\end{corollary}
\begin{proof}
On each connected component of $S$ the polarization gives one Hilbert
polynomial; apply \cref{akcp-6-3-picard-integral} and the decomposition in
\cref{akcp-5-8-spl-polynomial-open}.
\end{proof}

\begin{definition}[Relative dualizing sheaf]
\label{akcp-6-5-dualizing}
\dcref{6.5}
\source{MR0555258-compactifying-picard:page-0046}
\source{MR0555258-compactifying-picard:page-0047}
For a flat finitely presented proper family with Cohen--Macaulay pure
$r$-dimensional fibers there is a flat locally finitely presented sheaf
$\omega=\omega_{X/S}$ restricting to a dualizing sheaf on each fiber, with a
trace $R^rf_*\omega\to\OO_S$, unique up to unique isomorphism.  The Gorenstein
locus is the open retrocompact locus where $\omega(s)$ is invertible.  If
$X\subset\PP_S(E)$, then
\[
\omega=\Ext^{e-r}_{\PP_S(E)}(\OO_X,\OO_{\PP_S(E)}(-e-1)).
\]
For $S$-flat $\II$,
\[
\Ext_X^q(\II,\omega)=
\Ext_{\PP(E)}^{q+e-r}(\II,\OO(-e-1)),
\quad
\Ext_X^q(\II,\omega)=0\; (q>r-\max\depth\II(s)).
\]
If the fibers are geometrically integral, $\omega$ is relatively
pseudo-invertible.
\end{definition}
\begin{proof}
The dualizing sheaf and trace are obtained by relative duality; the displayed
Ext formula and base-change follow from the change-of-rings spectral sequence,
because the other local Ext groups vanish on Cohen--Macaulay fibers.  The
vanishing bound follows from the fiberwise duality theorem and
\cref{akcp-1-10-open-ext}.  Over a field, testing at closed points shows that
$\omega$ is Cohen--Macaulay of dimension $r$, hence torsion-free; reducedness
of an integral fiber makes its generic rank one.  These properties spread over
the base by openness.
\end{proof}

\begin{theorem}[Cohen--Macaulay Picard components]
\label{akcp-6-6-picard-cm}
\dcref{6.6}
\source{MR0555258-compactifying-picard:page-0047}
\uses{akcp-6-2-parameter,akcp-6-5-dualizing,akcp-5-13-picard-open}
If $X/S$ is flat projective with geometrically integral Cohen--Macaulay fibers
of fixed dimension and polarization polynomial, every polynomial component of
$\Pic^{\mathrm{CM},\acute{\mathrm et}}_{X/S}$ is strongly quasi-projective.
\end{theorem}
\begin{proof}
Use $\FF=\omega_{X/S}$ in \cref{akcp-6-2-parameter}.  The vanishing in
\cref{akcp-6-5-dualizing} puts the Cohen--Macaulay Picard functor inside that
parameter functor, and the condition is open retrocompact by
\cref{akcp-5-13-picard-open}.
\end{proof}

\begin{corollary}[Locally projective Cohen--Macaulay Picard]
\label{akcp-6-7-picard-cm-local}
\dcref{6.7}
\source{MR0555258-compactifying-picard:page-0047}
\uses{akcp-6-6-picard-cm,akcp-5-8-spl-polynomial-open}
For flat finitely presented locally projective $X/S$ with integral
Cohen--Macaulay fibers, the Cohen--Macaulay Picard functor is represented by a
separated locally finitely presented $S$-scheme; if $X/S$ is projective it is a
disjoint union of quasi-projective schemes.
\end{corollary}
\begin{proof}
The assertion is local on $S$.  In the projective case apply
\cref{akcp-6-6-picard-cm} on each connected component and use polynomial
decomposition.
\end{proof}
